% 本文件由 LaTeX 排版 Agent 根据有效 translation.json 自动生成。
% author=Ambrose; work_id=3411; title=论交出大殿

\chapter{论交出\OriginalTerm{大殿}{Basilicas}}

% structure: p0001 source=h1 target=omitted reason=page-title-already-rendered-by-parent

% structure: p0002 source=h2 target=omitted reason=duplicate-page-title

为了平息人民对皇帝法令的焦虑，他将自己的答复摆在他们面前，并补充说他之所以没有去\OriginalTerm{法庭}{consistory}，是因为害怕失去\OriginalTerm{大殿}{basilica}。随后，他首先向对手提出在圣堂内讨论，称自己并不畏惧他们的武器；并且，在重提他关于圣器的答复之后，宣告自己已准备好应战。他坚持说，天主的旨意不能被挫败，祂的保护也不能被胜过，但祂也预备在祂的仆人内受苦。既然他此前尚未被带走，显然异端者是无缘无故地制造这场骚乱。接着，在将\OriginalTerm{拿波特}{Naboth}的往事和基督进入\Place{耶路撒冷}的事迹应用到当前局势之后，他谴责\OriginalTerm{奥森提乌斯}{Auxentius}的残酷法律，回应\OriginalTerm{亚略派}{Arians}的反对意见，并声明他将乐于在民众面前讨论此事。他补充说，奥森提乌斯已受那些他选为法官的外教人定罪，正如他曾受\OriginalTerm{保禄}{Paul}和基督定罪一样。这个异端者忘记了前一年，那时他曾向\OriginalTerm{凯撒}{Cæsar}提出同样的申诉；而亚略派，由于煽动对基督仆人的恶意，比犹太人更为恶劣：因为教会不属于凯撒，而是彰显基督的肖像。最后，他就自己的答复和圣诗再补充几句，声明自己并非不服从，皇帝是教会的儿子，而奥森提乌斯比一个犹太人更坏。\par

1. 我见你们异常不安，并且密切关注着我。我想知道这是什么原因？是否因为你们看到或听说我从\OriginalTerm{护民官}{tribunes}手中接到了一项皇帝的命令，大意是我可以随意离开此地，而且所有愿意跟随我的人都可以随我而去？你们是否害怕我会抛弃教会，因担忧自身安全而离弃你们？但你们可以注意到我传的口信：抛弃教会的念头从未进入我的心中；因为我敬畏宇宙的天主，胜于敬畏地上的皇帝；如果有人要用武力将我拖出教会，我的身体固然可以被赶出，但我的心志却不能。我已准备好，如果他做出皇帝权力惯常之事，我就承受一位司铎必须承担的命运。\par

2. 那么，你们为什么不安呢？我绝不会自愿离弃你们，尽管若有人使用武力，我无法抵挡。我能够悲伤、哭泣、呻吟；面对武器、士兵、哥特人，我的眼泪就是我的武器，因为这是一位司铎的防御。我不应该、也不能以任何别的方式抵抗；但逃跑并抛弃教会不是我的作风；免得有人以为我这样做是出于畏惧某种更重的惩罚。你们自己知道，我习惯于尊敬我们的皇帝，但不向他们屈服，主动献身受罚，而不畏惧为我预备的一切。\par

3. 我多么希望我能确知教会永不会被交给异端者。如果这符合司铎的职分，我很乐意前往皇帝的宫殿，在宫殿而不是在圣堂里进行我们的讨论。但在\OriginalTerm{法庭}{consistory}上，基督通常不是被告，而是审判者。谁会否认信仰的事由应当在圣堂内辩护呢？如果有人有信心，就让他到这里来；不要寻求皇帝的裁判，那裁判已经显示出偏见，它通过已经通过的法律清楚地证明皇帝反对信仰；也不要寻求某些人的预期好感，他们想两边讨好。我不会以这样的方式行事，以至于给任何人提供通过对基督行不义而赚钱的机会。\par

4. 周围的士兵，环绕圣堂兵器的撞击声，并不摇动我的信德，但它们使我忧虑，因为我担心在把我留在这里时，你们可能会遭遇生命危险。因为如今我已学会了不畏惧，但我确实开始为你们更加担心。我恳求你们，允许你们的主教去迎战他的仇敌。我们有一个攻击我们的仇敌，因为我们的仇敌魔鬼，\BibleQuote{如同咆哮的狮子巡游，寻找可吞噬的人，} \BibleRef{伯前 5:8} 正如宗徒所说的。毫无疑问，他得到了——他得到了（我们并非受骗，而是对此早有警醒）这样试探的权能，以免我或许因身体的创伤而偏离信德的恳切。你们读过魔鬼如何用许多方式试探圣\OriginalTerm{约伯}{Job}，以及最后他如何求得并且得到了权能去试炼他的身体，使之长满毒疮。\par

5. 当有人提议我应该交出教会的圣器时，我作了如下答复：无论我的私产中的什么被索取，无论是田产、房屋、金子或银子——事实上，凡在我权下的任何东西，我都乐意交出。但我不能从天主的殿宇中取走任何东西；也不能交出我所领受用以保管而非予以交出的东西。我这样做是为皇帝的好处着想，因为无论是我交出或他接受，都是不合宜的。愿他聆听一位直言不讳的主教的话，如果他愿为自己做最好的事，就让他停止对基督行不义。\par

6. 这些话充满谦卑，且如我所想，充满了一位主教应当向皇帝展现的精神。但既然 \BibleQuote{我们的搏斗，不是对抗血和肉，而是}（更甚者）\BibleQuote{对抗天界的恶神，} \BibleRef{弗 6:12} 那试探者魔鬼便借着他的仆役使这战斗更加艰难，并想通过我肉身的创伤试探我。弟兄们，我知道我们为基督的缘故所承受的这些创伤，不是毁灭生命的创伤，反而是延展生命的创伤。请允许这竞赛进行吧。你们要做旁观者。反思一下，如果一座城拥有一名运动员，或一位精通某样其他高尚技艺的人，它便会急切地推举他出赛。在较小的事上你们惯于希望如此，为何在更重要的事上却拒绝这样做呢？不畏惧死亡、不被任何肉身的享乐所束缚的人，也不畏惧武器和蛮族。\par

7. 事实上，如果主已命定我参与这场斗争，那么你们如此多日多夜的不眠守候便是徒劳。基督的旨意必将实现。因为我们的主耶稣是全能的，这是我们的信德：因此祂旨意所要成就的事必将实现，而我们不可阻挠天主的旨意。\par

8. 你们今日听到了所读的：救主吩咐了宗徒们，把一匹驴驹牵来给祂，并嘱咐他们，若有人阻挡他们，就说：\BibleQuote{主要用它。}\BibleRef{路 19:35} 但现在，如果祂也同样地命令了那驴驹——也就是那惯于驮负重担的牲口的驹子，如同人必须如此，对人曾说过：\BibleQuote{你们凡劳苦负重担的都到我这里来，我要使你们得安息；你们背起我的轭吧，因为它是轻松的；}——我说，如果祂现在也命令了那驴驹被牵来给祂，差遣那些宗徒们，他们已经脱下肉身，显现为天使的形貌，是我们肉眼看不见的，这又如何呢？如果有人阻止他们，他们难道不会说：主要用他？例如，如果对现世生命的贪恋，或者血肉之情，或人间的交际（因为我们或许看起来使某些人喜爱）要阻止他们呢？而那在此世爱我的人，如果肯让我成为基督的牺牲，那将更显示他的爱，因为\BibleQuote{离开此世与基督同在，是好得无比的；但留在肉身内，为你们却是更为要紧。}\BibleRef{斐 1:23} 因此，亲爱的弟兄们，你们不要害怕。因为我知道，无论我遭受什么，我都是为基督而承受的。我也曾读到，我不该怕那只能杀害肉身的\BibleRef{玛 10:28}。我也曾听见那一位说过：\BibleQuote{谁若为我的缘故丧失自己的性命，必要获得性命。}\BibleRef{玛 10:39}\par

9. 因此，如果主愿意，肯定没有人能抵抗。如果祂仍然延缓我的斗争，你们害怕什么呢？守护基督仆人的，不是肉身的护卫，而是主的眷顾，惯常会保护他。\par

10. 你们因为发现双扇门开着而忧心，据说是有一个瞎子在寻找他住处时不慎打开的。从这事上你们可以明白，人的警觉并非保障。看！一个失去视力的人竟突破了我们的所有防御，躲过了守卫的注意。但主却从未丧失祂怜悯的守护。正如你们记得的，两天前不是也发现大殿左边某个入口开着，而你们原以为已经关好锁牢了吗？武装的人包围了大殿，试了这个入口又试那个，但他们的眼睛却瞎了，竟看不见那个开着的入口。你们也很清楚，那入口开了许多夜。因此，不要再焦虑了；因为基督所命令的、最有益的事必将发生。\par

11. 现在，我给你们举出法律书中的一些例子。厄里叟被叙利亚王搜寻；一支军队被派来捉拿他；他四面受困。他的仆人开始害怕，因为他是个仆人，也就是说，他没有自由的心灵，也没有自由的行动能力。那位圣先知设法开启他的眼睛，对他说：\BibleQuote{你观看，看看在我们这一边的比反对我们的更多。}\BibleRef{列下 6:16} 他就看见了，见到了成千的天使。因此，要记住，是那些看不见的，而非看得见的，守护着基督的仆人们。但如果他们守护你们，那是应允你们的祈祷：因为你们读过，那些搜寻厄里叟的人进入了撒玛黎雅，来到他们所想要捉拿的人那里。他们不但不能损害他，反而因着他们所敌对的人的代祷而得到了保护。\par

12. 宗徒伯多禄也给你们提供了这两种情况的一个例子。当黑落德搜寻他、捉住他时，他被投进了监狱。因为天主的仆人并没有逃走，而是坚定地站在那里，丝毫没有恐惧的念头。教会为他祈祷，但宗徒却在监狱中睡着了，这证明他并不害怕。一位天使被派来唤醒睡梦中的他，由天使引领伯多禄出了监狱，暂时免于一死。\par

13. 后来，伯多禄在击败西满后，在民众中播种天主的道理、教导贞洁，激起了外邦人的反感。当这些人搜寻他时，基督徒们请求他暂时躲避一下。虽然他渴望受苦，但他看到人们祈祷，便动了心，因为他们求他保全自己，以便教导和坚固他的子民。还需要我多说吗？夜间，他开始离开城镇，在城门口看见基督迎面而来，走进城里，便问：\Quote{主，祢往哪里去？}基督回答说：\Quote{我来是要再次被钉十字架。}伯多禄明白，这天主的答复是指向他自己的十字架，因为基督不可能再次被钉十字架，因为祂已经藉着所受的苦难死亡，脱去了肉躯；因为：\BibleQuote{祂死，是死于罪恶，仅有一次；祂活，是活于天主。}\BibleRef{罗 6:10} 就这样，伯多禄明白了基督要在祂的仆人身上再次被钉十字架。因此他甘愿返回；当基督徒们问他时，就告诉了他们原因。他随即被捕，并借着自己的十字架光荣了主耶稣。\par

14. 那么，你们看，基督愿意在祂的仆人身上受苦。如果祂对这个仆人说：\BibleQuote{我如果要他存留直到我来，你跟随我}，\BibleRef{若 21:22} 并愿意品尝这棵树的果实，那又如何呢？因为如果祂的食物是承行祂圣父的旨意\BibleRef{若 4:34}，那么，同样地，分受我们的苦难也是祂的食物。为了从主自己身上举个例子——祂难道不是愿意时才受苦，寻找时才被找到吗？但当祂受难的时辰尚未到来时，祂就从那些寻找祂的人中间穿过\BibleRef{若 7:30}，尽管他们看见了祂，却不能抓住祂。这清楚地告诉我们，当主愿意时，每个人就被找到并捉住；但因时候被推迟，即使他与寻找他的人面对面，也不会被捉住。\par

15. 难道我没有自己每天出去拜访，或前往殉道者的坟墓吗？我在往返时难道没有经过皇宫吗？然而，没有人向我动手，尽管他们有意驱逐我，正如他们后来散布的，说：\Quote{离开这城，随你去哪里吧。}我承认，我在寻求某种伟大的事，或是刀剑或是烈火，为\OriginalTerm{基督之名}{Name of Christ}；但他们却给我愉快的事，而非苦难；但\OriginalTerm{基督的竞技者}{Christ's athlete}不需要愉快的事，而需要苦难。因此，不要让任何人搅扰你们，因为他们提供了一辆马车，或是因为\OriginalTerm{\Person{奥克森提乌斯}}{Auxentius}——那个自称主教的人——说了些（他认为的）难听的话。\par

16. 许多人说，刺客已被派出，死刑已被判决于我。我毫不惧怕这一切，也不打算离弃我在此的位置。我能往哪里去呢？当没有一个心灵不充满呻吟和眼泪；当各处天主教会的主教们或被驱逐，或如果抵抗便遭刀剑，而每一个不服从命令的元老都遭取缔。而这些事都是由一位主教亲手所写、亲口所说，他为了显示自己最为博学，竟不忘怀一个古老的警告。因为我们在先知书中读到，他看见一把飞行的镰刀。\BibleRef{匝 5:1}\Person{奥克森提乌斯}效法此事，向各城发出飞行的刀剑。然而，\OriginalTerm{撒殚}{Satan}也化作\OriginalTerm{光明的天使}{angel of light}，\BibleRef{格后 11:14}并模仿其能力以行恶事。\par

17. 主耶稣，祢在一瞬间救赎了世界：难道\Person{奥克森提乌斯}能在一瞬间，尽其所能，杀戮如此多的民众，有的用刀剑，有的用亵渎？他带着血污的嘴唇和沾满鲜血的手，来索要我的\OriginalTerm{圣殿}{basilica}。今天宣读的经文正好回答他：\BibleQuote{但恶人，天主说：你怎敢传述我的正义？}这就是说，和平与疯狂之间没有联盟，基督与\OriginalTerm{贝里雅耳}{Belial}之间没有联盟。\BibleRef{格后 6:15}你们也记得今天我们读到\Person{纳波特}，一个拥有自己葡萄园的圣洁之人，在国王的要求下被逼交出葡萄园。当国王拔除葡萄树，打算种植普通蔬菜时，他回答他说：\BibleQuote{天主不容我交出我祖先的产业。}\BibleRef{列上 21:3}国王忧闷，因为正当地拒绝了他人的合法请求，却以女人的诡计不正当取得了。\Person{纳波特}用自己的血捍卫了他的葡萄树。如果他没有交出葡萄园，难道我们要交出基督的教会吗？\par

18. 难道我当时给出的答复是傲慢的吗？因为当我被传召时，我说：\Quote{天主不容我交出基督的产业。}如果\Person{纳波特}没有交出他祖先的产业，难道我要交出基督的产业吗？我还进一步说：\Quote{天主不容我交出我祖先的产业，就是\OriginalTerm{\Person{狄奥尼修斯}}{Dionysius}的产业，他为了信仰而被放逐，死于异地；\OriginalTerm{精修圣人}{Confessor}\OriginalTerm{\Person{欧斯托吉乌斯}}{Eustorgius}的产业，\OriginalTerm{\Person{密索克莱斯}}{Mysocles}的产业，以及昔日所有忠信主教们的产业。}我以主教当有的方式回答：\Quote{让皇帝按皇帝当有的方式行事。他必须夺走我的生命，而不是我的信仰。}\par

19. 但我要把它交给谁呢？今天福音中的读经应当教导我们，所要求的是什么，又是谁在要求。你们听到读经说，当基督\BibleRef{路 19:35}骑在驴驹上，孩子们大声呼喊，犹太人便恼怒。最后他们对主耶稣说，命令祂制止他们。祂回答说：\BibleQuote{如果这些人默不作声，石头就要呼喊了。}\BibleRef{路 19:40}然后，进入圣殿，祂赶出兑换银钱的人、桌子，以及在圣殿中卖鸽子的人。这段经文被读到，并非出于我的安排，而是偶然；但它恰好切合当前时局。对基督的赞美，永远是对不信者的鞭笞。如今当基督受赞美时，异端者就说煽动骚乱的话。异端者说死亡正为他们预备，而确实，他们的死亡就在于对基督的赞美。他们怎能忍受对祂的赞美呢？他们可是主张祂的软弱啊。因此今天，当基督受赞美时，\OriginalTerm{亚略派}{Arians}的疯狂正遭鞭笞。\par

20. \OriginalTerm{\Place{革辣撒人}}{Gerasenes}不能容忍基督的临在；\BibleRef{路 8:37}这些人，比\Place{革辣撒人}更甚，竟不能忍受对基督的赞美。他们看见男孩们歌唱基督的荣耀，因为经上记载：\BibleQuote{从婴孩和幼童的口中，祢成就了赞美。}他们嘲笑他们的稚龄，那么充满信德，并说：\Quote{看，他们为什么呼喊？}但基督回答他们说：\BibleQuote{如果这些人默不作声，石头就要呼喊了，}\BibleRef{路 19:40}这就是说，更有力者将要呼喊，青年和年长将要呼喊，老年人将要呼喊；这些石头如今稳固地建在那磐石上，关于这磐石经上记载：\BibleQuote{匠人所弃的石头，反成了角石。}\par

21. 于是，受到这些赞美的邀请，基督进入祂的圣殿，\BibleRef{若 2:15}拿起鞭子，将兑换银钱的人赶出圣殿。因为祂不允许金钱的奴仆留在祂的圣殿里，也不允许那些售卖座位的人在那里。座位岂不是荣衔？鸽子岂不是那些追随纯洁、清晰信仰的纯朴心灵和灵魂？那么，我岂能把基督拒之门外的人领进圣殿？因为那些出卖尊荣和地位的人，必被吩咐离开。凡想出售信徒纯朴心灵的人，必被吩咐离开。\par

\newline{}22. 因此，\Person{奥森提乌斯}被驱逐。\Person{墨丘利乌斯}被拒之门外。同一怪物，两个名字！为了不让人知道他是谁，他改了名，自称\Person{奥森提乌斯}，因为这里曾有一位\OriginalTerm{亚略派}{Arian}主教，名叫\Person{奥森提乌斯}。他这样做是为了欺骗那先前在\Person{奥森提乌斯}权下的人民。他改了名，却没有改变他的虚伪。他脱下狼皮，却又穿上狼皮。他改名对他毫无帮助，无论如何他都为人所知。他在\Place{斯基提亚}被称呼一个名字，在这里被称呼另一个名字。他在每个居住的国家都有一个名字。他已经有两个名字了，如果他从这里去别处，还会有第三个名字。因为，他怎能忍受让一个名字成为如此邪恶的证据呢？他在\Place{斯基提亚}作的恶较小，便羞愧得改了名。在这里他竟敢作更大的恶，难道无论到哪里，他愿意让他的名字出卖他吗？他亲手签署了许多人的死刑判决，却还能心神不动摇吗？\newline{}\par

\newline{}23. 主耶稣只将少数人赶出祂的圣殿，但\Person{奥森提乌斯}却一个不留。耶稣用鞭子将他们赶出圣殿，\Person{奥森提乌斯}却用剑；耶稣用鞭子，\Person{墨丘利乌斯}用斧子。圣善的主用鞭子驱逐亵圣者；邪恶的人却用剑迫害圣人。关于他，你们今天说得好：让他带着他的法律走吧。他会带走的，即使他不愿意；他会带走他的良心，尽管他不带走任何文书；他会带走用鲜血铭刻的灵魂，尽管他不会带走用墨水写成的信。经上记载：\BibleQuote{\Person{犹大}，你的罪恶是用铁笔和金刚石的尖端记录的，铭刻在你的心版上，}\BibleRef{耶 17:1}就是说，它写在那里，它就是从那里发出来的。\newline{}\par

\newline{}24. 他，一个满手鲜血、满身杀戮的人，竟敢向我提及辩论？他认为那些不能用言语误导的人就当用剑来杀害，他用口颁布血腥的法律，用手写下它们，并以为法律能为人们规定应持有的信仰。他没有听过今天所读的：\BibleQuote{人成义不是靠法律的行为，}\BibleRef{迦 2:16}或\BibleQuote{我因法律已死于法律，为能生活于天主，}\BibleRef{迦 2:19}就是说，通过属神的法律，他对法律属肉性的解释已死。而我们，藉着我们的主耶稣基督的法律，已死于这批准如此背信指令的法律。法律并没有聚集起教会，而是基督的信德聚集的。因为法律不是本于信德，而是\BibleQuote{义人因信德而生活。}\BibleRef{迦 3:11}所以，使人成义的是信德，不是法律，因为义德不是通过法律，而是通过基督的信德。那抛弃自己的信德、为法律的规条辩护的人，正好证明自己是不义的；因为义人因信德而生活。\newline{}\par

\newline{}25. 那么，有人该随从这法律吗？依此法律，\OriginalTerm{阿里米农会议}{Council of Ariminum}得到批准，在会上基督被称为受造物。但他们说：\BibleQuote{天主派遣了祂的儿子，由女人所生，生于法律之下。}\BibleRef{迦 4:4} 所以他们如此说\Emph{被造的}，就是说\Emph{受造的}。难道他们没有考虑他们所引用的这些话：基督是\Emph{被造的}，却是由女人所生；就是说，就祂由童贞女诞生而言，祂是\Emph{被造的}，但就祂的天主出生而言，祂由父所生。他们今天也读过：\BibleQuote{基督由法律的咒骂中救赎了我们，为我们成了可咒骂的}\BibleRef{迦 3:13}吗？基督在祂的\OriginalTerm{天主性}{Godhead}上是可咒骂的吗？但是为什么祂被称为可咒骂的，宗徒告诉我们，记载说：\BibleQuote{凡被悬在木架上的，是可咒骂的，}\BibleRef{迦 3:13}也就是说，那在祂的肉身中承担了我们的肉身，在祂的身体中承担了我们的软弱和咒骂，为能将它们钉死的那一位；因为祂本身不是可咒骂的，而是在你们内成为可咒骂的。所以别处记载说：\BibleQuote{没有认识罪的，却替我们成了罪，因为祂承担了我们的罪过}\BibleRef{格后 5:21}为能藉着祂\OriginalTerm{受难的圣事}{Sacrament of His Passion}消灭它们。\newline{}\par

\newline{}26. 我的弟兄们，这些事，我本想当着你们的面与他详细讨论；但他知道你们并非不熟悉信德，便逃避在你们面前的审判，而选择了四五个外教人代表他，如果那是他确实选了人，我倒希望这些人在我们中间，不是为审判有关基督的事，而是为聆听基督的威严。然而，他们已经对\Person{奥森提乌斯}作出了裁决，当他日复一日在他们面前申辩时，他们都不相信他。还有什么比这更能定他的罪呢：在没有对手的情况下，竟在自己的法官面前被打败？所以，现在我们也握有他们反对\Person{奥森提乌斯}的意见。\newline{}\par

\newline{}27. 他选择了外教人，这是理当受谴责的；因为他无视宗徒的命令，宗徒说：\BibleQuote{你们中间有人与另一人有了争讼，怎么敢在不义的人面前起诉，而不在圣徒面前呢？你们不知道圣徒要审判世界吗？}\BibleRef{格前 6:1} 并在下面说：\BibleQuote{难道你们中间没有一个智慧人，能在外教人之间判断吗？但兄弟与兄弟互相控告，且是在不信主的人面前。}\BibleRef{格前 6:5} 所以，你们看，他所引入的是违反宗徒权威的。那么，你们决定吧，我们该追随\Person{奥森提乌斯}，还是以\Person{保禄}为师。\newline{}\par

28. 但为什么要提宗徒呢，当主亲自藉先知呼喊说：\BibleQuote{我的百姓，你们这些认识判断的，在我心上有我法律的，请听我！} \BibleRef{依 51:7} 天主说：\BibleQuote{我的百姓，你们认识判断的，请听我。} \Person{奥森提乌斯}说：\Quote{你们不认识判断。} 你看他怎样在你内定天主的罪，他拒绝天上神谕的声音，那声音说：\BibleQuote{我的百姓，请听我。} 主说。他没有说\Quote{外邦人，请听我}，也没有说\Quote{犹太人，请听我}。因为那些曾是主百姓的人，如今成了错误的百姓；而那些曾是错误百姓的人，已开始成为天主的百姓，因为他们信仰了基督。所以那百姓判断，在他们心中有神圣的，而非人的法律，那法律不是用墨水写成的，而是写在生活天主的神内；\BibleRef{格后 3:3} 不是记在纸上，而是印在心上。那么，是谁对你做错了呢？是那个拒绝被你聆听的人，还是那个选择被你聆听的人？\par

29. 他四面受困，便求助于他祖辈的诡计。他想要煽动对皇帝的恶意，说一个青年，一个对圣经无知的望教者，竟该在\OriginalTerm{御前会议}{consistory}中定案。仿佛去年我被召前往宫殿，当在众首领面前于御前会议讨论此事，皇帝想要夺取大堂时，我在当时宫廷面前被吓住了，没有表现出一位主教应有的坚定，或是让步了。难道他们不记得，当百姓得知我去了宫殿，便蜂拥而来，势不可挡；并且所有人都为基督的信仰甘愿赴死，因有军官带轻装部队出来驱散人群？我不是被请来以长篇演讲安抚百姓吗？我不是保证没有人会侵犯教会的大堂吗？虽然我被要求做一件善事，但百姓来宫殿这件事却用来攻击我。他们现在又想逼我再次如此。\par

30. 我召回了百姓，然而并未逃脱他们的恶意；这种恶意，我认为我们更应无畏面对而非惧怕。因为为什么我们要为基督之名而害怕呢？除非或许我应受扰于他们的话：\Quote{皇帝难道不应该有一座大堂可以去，而\Person{安博罗削}竟想比皇帝更有权力，因而拒绝皇帝前往教堂的机会？} 当他们这样说时，他们想要抓住我的话，如同犹太人用诡诈的话试探基督说：\BibleQuote{师傅，给\Person{凯撒}纳税，可以不可以？} \BibleRef{玛 22:17} 难道总是因凯撒的缘故煽动对天主仆人的恶意，而不敬虔者利用此来散布诽谤，以便隐藏在皇帝的名义下吗？他们能说自己没有分享那些其所从之人的亵渎吗？\par

31. 看，\OriginalTerm{亚略派}{Arians}比犹太人还坏得多。他们问祂是否认为应向凯撒纳税；这些人却想把教会权交给凯撒。但正如这些无信德者跟随他们的创始者，我们也该如同我们的主和创始者所教的那样回答。因为耶稣看透犹太人的恶意，对他们说：你们为什么试探我？拿一个税币给我看。当他们拿给他，他说：\BibleQuote{这肖像和名号是谁的？} \BibleRef{玛 22:18} 他们回答说：凯撒的。耶稣便对他们说：\BibleQuote{那么，凯撒的，就应归还凯撒；天主的，就应归还天主。} \BibleRef{玛 22:21} 我也这样对那些反对我的人说：给我看一个税币。耶稣看见凯撒的税币说：凯撒的，就应归还凯撒；天主的，就应归还天主。他们夺取教会的大堂时，能拿出凯撒的税币吗？\par

32. 可是在教会内，我只认识一个\OriginalTerm{肖像}{Image}，就是那无形天主的肖像，关于祂，天主曾说：\BibleQuote{让我们照我们的肖像，按我们的模样造人；} \BibleRef{创 1:26} 论到那肖像，经上记载：基督是祂荣耀的光辉，是祂本体的肖像。 \BibleRef{希 1:3} 在那肖像内我看见了父，正如主耶稣亲口说的：\BibleQuote{谁看见了我，就是看见了父。} \BibleRef{若 14:9} 因为这肖像并不与父分离，祂确曾教导了我圣三的合一，说：\BibleQuote{我与父原是一体，} \BibleRef{若 10:30} 又说：\BibleQuote{父所有的一切，都是我的。} \BibleRef{若 16:15} 关于圣神，也说：圣神是基督的，并由基督领受，正如经上所载：\BibleQuote{祂要由我领受，而传告给你们。} \BibleRef{若 16:14}\par

33. 那么，我们怎能回答得不够谦卑呢？如果他索要赋税，我们并不拒绝。教会的土地缴纳税。如果皇帝想要土地，他有权索要，我们没有人会干涉。百姓的捐献足够用于穷人。不要在土地的事上煽动恶意。他们若要，就拿去吧，若是皇帝的旨意。我不给予，但也不拒绝。他们索要金子。我可以说：我既不求金银。但他们却因金钱的花费而煽动恶意。我不怕这种恶意。我有受我赡养的人。我的赡养对象是基督的穷人。我知道如何收集这种财富。在这一点上，他们甚至可以控告我，说我花掉了穷人的钱！如果他们控告我向他们寻求保护，我不否认；不，我恳求这种保护。我有我的保护，那就是穷人的祈祷。瞎子瘸子，弱者老者，比强壮的战士更有力。总之，对穷人的赠予使天主欠我们的债，因为经上记载：\BibleQuote{谁施恩于穷人，便是贷款于上主。} \BibleRef{箴 19:17} 战士的护卫往往不配得天主的恩宠。\par

34. 他们还宣称，人们被我圣诗的曲调引入了歧途。我当然不否认这一点。那是一种崇高的曲调，没有什么比它更有力量。因为有什么比全体人民每日用口所庆祝的\OriginalTerm{天主圣三}{Trinity}的\OriginalTerm{宣认}{confession}更有力量呢？所有人都争先恐后地宣认\OriginalTerm{信德}{faith}，并且知道如何用诗歌赞美圣父、圣子及圣神。于是他们都成了老师，而他们几乎还不能做门徒。\par

35. 还有什么比我们效法基督的榜样更能表现出\OriginalTerm{服从}{obedience}呢？\Quote{祂贬抑自己，听命至死？}\BibleRef{斐 2:7} 因此祂藉着自己的服从解救了一切。\Quote{正如因一人的悖逆，大众都成了罪人；同样，因一人的服从，大众都成了义人。}\BibleRef{罗 5:19} 那么，如果祂是服从的，就让他们接受服从的法则吧：我们持守这法则，对那些在皇帝一方挑拨离间的人说：\Person{凯撒}的归\Person{凯撒}，天主的归天主。应向\Person{凯撒}纳税，我们并不否认。教会属于天主，所以它不应该被转让给\Person{凯撒}。因为天主的殿绝不能依法归于\Person{凯撒}。\par

36. 这话是出于对皇帝尊敬的感情，无人能否认。因为还有什么比皇帝被称为教会的儿子更充满尊敬呢？既然如此说，也是无罪而说的，因为这说法带着\OriginalTerm{天主的恩宠}{divine favour}。因为皇帝是在教会之内，而非在其上。一位好皇帝寻求教会的帮助，也不拒绝它。我既以全然的谦卑说这话，也以坚定加以申明。有人用火、剑、充军来威胁我们；我们作为基督的仆人，已经学会了不害怕。对于那些不害怕的人，没有什么事能成为真正可怕的缘由。因此经上也记着：\Quote{他们的打击，不过如同幼儿的箭矢。}\par

37. 那么，看来对他们的建议已作了充分的回答。现在我要问他们，救主所问过的话：\Quote{若翰的洗礼是从天上来的，还是从人来的？}\BibleRef{路 20:4} 犹太人无法回答祂。如果犹太人没有轻视若翰的洗礼，那么\Person{奥森提乌斯}难道轻视基督的\OriginalTerm{圣洗}{baptism}吗？因为那圣洗不是人间的洗礼，而是从天上来的，是\OriginalTerm{伟大谋士的使者}{angel of great counsel}\BibleRef{依 9:6}带给我们的，好使我们能在天主前\OriginalTerm{成义}{justification}。那么，为什么\Person{奥森提乌斯}认为，已经因\OriginalTerm{天主圣三}{Trinity}之名受洗的信友应当重新受洗呢？宗徒明明说过：\Quote{只有一个\OriginalTerm{信德}{faith}，一个\OriginalTerm{圣洗}{baptism}}？\BibleRef{弗 4:5} 他既然藐视天主的计划，并谴责基督为救赎我们的罪而赐予我们的圣洗，又怎能说自己不是基督的仇敌，而只是人的仇敌呢？\par

\SourceRecord{Translated by H. de Romestin, E. de Romestin and H.T.F. Duckworth.；Nicene and Post-Nicene Fathers, Second Series；Vol. 10.；Edited by Philip Schaff and Henry Wace.；Buffalo, NY: Christian Literature Publishing Co.,；1896.；<http://www.newadvent.org/fathers/3411.htm>.}
